% Part: applied-modal-logics
% Chapter: temporal-logic
% Section: possible-histories

\documentclass[../../../include/open-logic-section]{subfiles}

\begin{document}

\olfileid[tr]{aml}{tl}{poss}

\olsection{Olası Tarihçeler}

Erişilebilirlik bağıntısına herhangi bir kısıtlama getirmemiz gerekmediğinden,
şimdiye dek kullandığımız zamansal mantığın bağıntısal modelleri son derece
esnektir. Bu, zamansal modellerin hem geçmişte hem gelecekte dallanabileceği
anlamına gelir. Fakat olay dizilerini olası tarihçeler olarak ele aldığımız,
dallanmaya ilişkin daha ``modal'' bir kavrayışı incelemek isteyebiliriz. Bunun
için dilimizi değiştirmek zorunda değiliz; yine de ``olağan'' $\Box$ ve
$\Diamond$ modal işleçlerimizi ekleyebilir ve öznelerin bilgisinin zaman
içindeki değişimini temsil etmek için epistemik erişilebilirlik bağıntıları
eklemeyi de düşünebiliriz.

\begin{defn}
  Zamansal dil için bir \emph{olası tarihçeler modeli},
  $\mModel{M}=\tuple{T,C,V}$ üçlüsüdür; burada
  \begin{enumerate}
    \item $T$, zaman içindeki durumlar olarak yorumlanan boş olmayan bir kümedir.
    \item $C$, hesaplama yollarından ya da bir sistemin \emph{olası
      tarihçelerinden} oluşan bir kümedir. Başka deyişle $C$, $\sigma$ ile
      gösterilen ve $s_1$, $s_2$,~$s_3$, \dots durumlarından oluşan diziler
      kümesidir; burada her $s_i\in T$'dir.
    \item $V$, her !!{propositional variable}~$p$'ye zaman noktalarından
      oluşan bir $V(p)$ kümesi atayan bir fonksiyondur.
  \end{enumerate}
  Konuyu basitleştirmek için genellikle, bir tarihçe $C$ içindeyse onun bütün
  son parçalarının da orada bulunduğunu varsayacağız. Örneğin $s_1$,
  $s_2$,~$s_3$ dizisi $C$ içindeyse $s_2$, $s_3$ dizisi ve ayrıca $s_3$ de
  öyledir. Ayrıca iki $s_i$ ve~$s_j$ durumu bir $\sigma$ dizisinde
  yer aldığında, $s_i\prec_\sigma s_j$ ifadesini $i<j$ olduğunda kullanırız.
  $t\in V(p)$
  olduğunda $p$'nin $t$'de \emph{doğru} olduğunu söyleriz.
\end{defn}

İlgili tek değişiklik, bir $t$ zaman noktasındaki !!a{formula} doğruluğunu
bir $\mModel M$ modelinde değerlendirirken bunu bir $\sigma$ tarihçesine
göre yapmamızdır; bu tarihçede $t$ bir durum olarak yer alır. Buna karşılık
!!{propositional variable}s ya da doğruluk fonksiyonlu bağlaçların hiçbirinin
semantiğini değiştirmemiz gerekmez. Bunların tümü tam olarak
\olref[sem]{defn:tmodels}'deki gibidir; çünkü hiçbiri $\sigma$'ya gönderimde
bulunmaz. Fakat şimdi gelecek işlecimiz $\Ftemp$'i yeniden tanımlıyor ve bu
tarihçelere göre $\Diamond$ işlecimizi ekliyoruz.

\begin{defn}\ollabel{defn:phmodels}
  \emph{!!a{formula}~$!A$'nın $t,\sigma$'daki doğruluğu}, $\mModel M$ içinde,
  yani $\mSat{M}{!A}[t,\sigma]$:
  \begin{enumerate}
  \item\ollabel{defn:sub:phmodels-f} \indcase{!A}{\Ftemp
    !B}{$\mSat{M}{\indfrm}[t,\sigma]$ ancak ve ancak
    $\mSat{M}{!B}[t',\sigma]$ olan bazı $t'\in T$ bulunup
    $t\prec_\sigma t'$ ise geçerlidir.}
  \item\ollabel{defn:sub:phmodels-diamond} \indcase{!A}{\Diamond
    !B}{$\mSat{M}{\indfrm}[t,\sigma]$ ancak ve ancak
    $\mSat{M}{!B}[t,\sigma']$ olan bazı $\sigma'\in C$ bulunur ve
    $t$ bu tarihçede yer alırsa geçerlidir.}
  \end{enumerate} 
\end{defn}

Diğer zamansal ve modal işleçler benzer biçimde tanımlanabilir. Ancak artık
zaman kipini ve modaliteyi birleştiren iddiaları temsil edebiliriz. Örneğin
``$p$ gerçekleşmeyecek ama gerçekleşebilirdi'' iddiasını
$\lnot\Ftemp p\land\Diamond\Ftemp p$ formülüyle simgeleyebiliriz. Bu formül,
$p$'nin bir ardıl durumda doğru hâle gelmediği bir nokta ve tarihçede,
$p$'nin doğru hâle geleceği alternatif bir tarihçe bulunduğunda geçerli olur.

\end{document}
